\begin{corollary}\label{cor:smooth-dual-exponential-limit}
Under the hypotheses of \cref{cor:smooth-dual-exact-equation}, let $u$ be the threshold solution with $\sup_Xu=0$. If $u_\lambda$ is the unique bounded solution of
\[
 \mathrm{MA}_\omega(u_\lambda)=\mathrm{e}^{\lambda u_\lambda}f\,\mathrm{d}V,
 \qquad\lambda>0,
\]
then $u_\lambda-\sup_Xu_\lambda\to u$ uniformly as $\lambda\downarrow0$.
\end{corollary}
\begin{proof}
By \cref{cor:smooth-dual-exact-equation}, $f$ agrees almost everywhere with a smooth function, so it belongs to $L^p$ for every finite $p$. By \cref{prop:individual-normalizer-lp}, the constant $c$ for which $\mathrm{MA}_\omega(v)=cf\,\mathrm{d}V$ has a bounded solution is $b_\omega(f)$. Thus \cref{prop:symmetric-equation-clusters}, with the trivial symmetry group, shows that the normalized family $u_\lambda-\sup_Xu_\lambda$ is uniformly precompact as $\lambda\downarrow0$ and that every cluster limit is a normalized solution with this constant. Each cluster limit is therefore a normalized threshold subsolution and equals $u$ by \cref{cor:smooth-dual-exact-equation}. If the whole family did not converge uniformly, a sequence staying at a fixed positive uniform distance from $u$ would have a uniformly convergent subsequence, a contradiction.
\end{proof}

\begin{corollary}\label{cor:kahler-smooth-dual-attainment}
Suppose $\omega$ is K\"ahler. For a nonzero nonnegative $f\in L^1(\mathrm{d}V)$, the functional $D_{f^{1/n}}$ has a smooth strictly positive maximizer in $\mathscr G_\omega$ if and only if $f$ agrees almost everywhere with a smooth strictly positive function.
\end{corollary}
\begin{proof}
If a smooth strictly positive maximizer exists, the proof of \cref{cor:smooth-dual-exact-equation} gives
\[
 \omega+\ddc u=a\gamma,
 \qquad a\geq0,
\]
with $a,u$ smooth and $a\not\equiv0$. Since $\mathrm{d}\omega=0$, the first-order equation in \cref{lem:smooth-conformal-current} becomes $\mathrm{d}a=-aC$ for a smooth one-form $C$. Along any smooth path $\sigma$, this is the linear scalar equation
\[
 (a\circ\sigma)'=-(a\circ\sigma)C(\dot\sigma),
\]
whose solution is the initial value multiplied by a positive exponential. Since $X$ is connected and $a\not\equiv0$, it follows that $a>0$ everywhere, and $f=b_\omega(f)^{-1}a^n\gamma^n/\mathrm{d}V$ is smooth and strictly positive. Conversely, if $f$ is smooth and strictly positive, \cref{thm:weak-gaud-compact-duality}, applied to $F=\log f$, for which $\mathfrak D_F=D_{f^{1/n}}$, provides a smooth strictly positive maximizer.
\end{proof}

For nonclosed backgrounds, the strict positivity of $f$ in \cref{cor:kahler-smooth-dual-attainment} can fail on every compact complex manifold of dimension at least two, including K\"ahler manifolds.

\begin{example}\label{ex:nonclosed-dual-zero-density}
On every connected compact complex manifold $X$ of dimension $n\geq2$, with a fixed smooth strictly positive probability volume $\mathrm{d}V$, there is a nonclosed Hermitian form $\omega$ and a nonzero smooth nonnegative density $f$ vanishing on an open set such that $b_\omega(f)=1$ and $D_{f^{1/n}}$ has a smooth strictly positive maximizer.
\end{example}
\begin{proof}
Fix a Gauduchon metric $\gamma$ and a point $x$. Choose a smooth real function $\psi$ such that $\ddc\psi\leq-\gamma$ on a coordinate neighborhood $U$ of $x$; in holomorphic coordinates $z=(z_1,\ldots,z_n)$ centered at $x$, such a function is obtained by multiplying $-M|z|^2$, with $M$ large enough, by a cutoff that equals one near $x$, and then shrinking $U$.

Choose a smooth function $\chi$ with $0\leq\chi\leq1$ that vanishes on a smaller neighborhood of $x$ and equals one outside $U$. Put
\[
 a=L\chi,\qquad \omega=a\gamma-\ddc\psi.
\]
Since $X$ is compact, $\omega$ is strictly positive for sufficiently large $L>0$: on $U$ the term $-\ddc\psi$ is at least $\gamma$, and outside $U$ the term $L\gamma$ dominates $\ddc\psi$. Set $f=a^n\gamma^n/\mathrm{d}V$. The smooth function $\psi$ satisfies $\omega+\ddc\psi=a\gamma\geq0$ and $\mathrm{MA}_\omega(\psi)=f\,\mathrm{d}V$, so $b_\omega(f)\geq1$.

Stokes' theorem and $\ddc\gamma^{n-1}=0$ give
\[
 \int_X\omega\wedge\gamma^{n-1}
 =\int_Xa\gamma^n
 =\int_X(f\,\mathrm{d}V/\gamma^n)^{1/n}\gamma^n.
\]
Testing \eqref{eq:integrable-gauduchon-threshold} with $\gamma$ therefore gives the opposite inequality $b_\omega(f)\leq1$. Hence $b_\omega(f)=1$, the metric $\gamma$ attains the infimum, and the normalized current $\gamma^{n-1}/\int_X\omega\wedge\gamma^{n-1}$ is a smooth strictly positive maximizer of $D_{f^{1/n}}$. If $\omega$ were closed, \cref{cor:kahler-smooth-dual-attainment} would force $f$ to be strictly positive, contrary to its vanishing on an open set. Thus $\omega$ is nonclosed.
\end{proof}

In \cref{ex:nonclosed-dual-zero-density}, \cref{cor:smooth-dual-exact-equation,cor:smooth-dual-exponential-limit} show that every threshold subsolution is a constant translate of $\psi$ and that the normalized solutions of the exponential equations converge uniformly to $\psi-\sup_X\psi$. Thus vanishing of the density on an open set does not by itself force nonuniqueness or the failure of this approximation.

The restriction $n\geq2$ in \cref{cor:smooth-dual-exact-equation} is necessary. On a compact connected Riemann surface $X$ with a Hermitian form $\omega$, a normalized Gauduchon current has bidegree $(0,0)$, and the equation $\ddc T=0$ makes it a harmonic function, hence the constant $(\int_X\omega)^{-1}$. This single smooth strictly positive current maximizes for every nonzero nonnegative $L^1$ density, which need not have a smooth representative.
